% 20260522_Assingment_Probability_Statistics
\documentclass[dvipdfmx]{jsreport}
\usepackage{soupreamble}


\title{確率及び統計学特論A 第2回課題}
\author{M1 6012 川村聡一郎}
\date{}

\begin{document}
\maketitle

\begin{tcolorbox}
	確率変数$X, Y$に対して
\begin{equation*}
	\calP(X\le a)= \calP(Y\le a),\ a\in \R
\end{equation*}
	ならば，
\begin{equation*}
	\calP(X\in A)= \calP(Y\in A),\ A\in \calB(\R)
\end{equation*}
	であることを示せ．
\end{tcolorbox}


\begin{souanswer} %答案はここに書く
	実数$a\in \R$に対して，区間$I_a= (-\infty, a]$を考える．このような区間全体からなる集合族を$\calP$とおく．
\begin{equation}
	\calP= \{(-\infty, a]\; a\in \R\}
\end{equation}
	$a, b\in\R$に対して$(-\infty, a]\cap (-\infty, b]= (-\infty, \min{(a, b)}]\in \calP$であるためには$\calP$は共通部分について閉じている．
	すなわち$\calP$は$\pi$系であり，Borel集合族の定義より
\begin{equation}
	\sigma(\calP)= \calB(\R)
\end{equation}
	確率が一致するようなBorel集合の集合を$calL$とする．
\begin{equation}
	\calL= \{A\in \calB(\R); \calP(X\in A)= \calP(Y\in A)\}
\end{equation}
	この$\calL$が$\lambda$系の条件を満たすことを示す．
\begin{enumerate_alpha}
	\item $\R\in \calL$：
	$\calP(X\in A)= 1$かつ$\calP(Y\in A)= 1$であるから，
\begin{equation}
	\calP(X\in A)= \calP(Y\in A)
\end{equation}
	となり，$\R\in \calL$が成り立つ．

	\item $A, B\in \calL, A\subset B\Rightarrow B\setminus A\in \calL$：
	$A, B\in \calL, A\subset B$とする．確率の加法性より
\begin{align}
	\calP(X\in B\setminus A)&= \calP(X\in B)- \calP(X\in A)\\
	\calP(Y\in B\setminus A)&= \calP(Y\in B)- \calP(Y\in A)
\end{align}
	$A, B\in \calL$の定義より，$\calP(X\in A)= \calP(Y\in A)$かつ$\calP(Y\in B)= \calP(Y\in B)$であるから
\begin{equation}
	\calP(X\in B\setminus A)= \calP(Y\in B\setminus A)
\end{equation}
	であり，$B\setminus A\in \calL$
	\item $\{A_n\}\in \calL$が単調増加ならば$\bigcup_{n= 1}^\infty A_n\in \calL$：
	互いに素な可算個のBorel集合列$\{A_n\}\subset \calL$をとる．確率の可算加法性より
\begin{align}
	\calP\left(X\in \bigcup_{n= 1}^\infty A_n\right)&= \sum_{n= 1}^\infty \calP(X\in A_n)\\
	\calP\left(Y\in \bigcup_{n= 1}^\infty A_n\right)&= \sum_{n= 1}^\infty \calP(Y\in A_n)
\end{align}
	各$n$について$A_n\in \calL$であり，$\calP(X\in A_n)= \calP(Y\in A_n)$であるから各項がすべて等しく，無限級数の和も等しくなる．
\begin{equation}
	\calP\left(X\in \bigcup_{n= 1}^\infty A_n\right)= \calP\left(Y\in \bigcup_{n= 1}^\infty A_n\right)
\end{equation}
	となり，$\bigcup_{n= 1}^\infty A_n\in \calL$
	よって$\calL$は$\lambda$系
\end{enumerate_alpha}
	任意の$a\in \R$に対して$\calP(X\le a)= \calP(Y\le a)$が成り立つ．これは$\calP\subset \calL$を意味する．ここで$\calP$は$\pi$系，$\calL$は$\lambda$系である．$\pi\mathrm{-}\lambda$定理より
\begin{equation}
	\sigma(\calP)\subset \calL
\end{equation}
	が成り立ち，$\sigma(\P)= \calB(\R)$であるから
\begin{equation}
	\calB(\R)\subset \calL
\end{equation}
	が得られ，$\calL$の定義より任意の$A\in \calB(\R)$に対して$\calP(X\le a)= \calP(Y\le a)$が成立．
\end{souanswer}
\end{document}
